% 結論と今後の展望

% ----- 結論 -----

本論文では、日本のアニメ産業におけるフリーランスアニメーターの作業の体系的な不可視性に対処する
ブロックチェーンベースの貢献証明フレームワーク\textbf{AnimatorSBT}を提示した。
既存の制作管理ツールとの統合を前提としたSoulbound Tokenの発行を設計することで、
AnimatorSBTは制作タスクの粒度での個人の貢献について、
不変で検証可能かつポータブルな記録の作成を目指すものである。

我々の分析により、明確な研究ギャップが特定された：
SBTは医療、教育、農業に応用されており、ブロックチェーンはクリエイティブIPの管理にも活用されているが、
アニメ制作ワークフローにおける個人レベルの貢献証明を対象とした先行研究は存在しない。
我々は貢献の不可視性問題を形式化し、
理想的な解決策のための5つの要件（完全性、検証可能性、永続性、ポータビリティ、最小限の摩擦）を定義し、
AnimatorSBTがその5つすべてを充足することを示した。

プロトタイプ実装---Polygon Amoyテストネットに展開したスマートコントラクト、
発行サービス（Issuance Service）、
および検証ポータル（Verification Portal）---により、
実測ガスコストに基づく設計の検証を行った。
これらの実測値を用いたコストシミュレーションは本アプローチの経済的実現可能性を示す：
単一の貢献証明書のミントコストは約\$0.0009であり、
業界全体での完全導入（日本のフリーランスアニメーター全員）でも年間わずか\$256---
業界の年間収益の$0.000011\%$に過ぎない。
これにより、ブロックチェーンベースの認証に対する技術的・経済的障壁が
Layer~2スケーリングソリューションによって実質的に解消されたことが確認された。

% ----- 今後の展望 -----

本研究からいくつかの今後の研究方向が浮かび上がる：

\begin{enumerate}
  \item \textbf{ユーザースタディ：}
    現役のアニメーターおよび制作管理者を対象とした正式なユーザビリティ評価を実施し、
    認知的価値、信頼性、導入意欲を評価する。
    本研究にはフリーランサー（SBT受領者）と
    スタジオ（潜在的な検証者・共同署名者）の双方を含めるべきである。

  \item \textbf{スタジオ共同署名プロトコル：}
    スタジオがアニメーターの貢献主張を承認する二者間確認メカニズムを設計・実装し、
    自己証明のSBTを相互検証された認証情報へと変換する。
    これにはスタジオの参加動機を確保するための慎重なインセンティブ設計が必要である。

  \item \textbf{W3C Verifiable Credentialsとの相互運用性：}
    AnimatorSBTをW3C Verifiable Credentials Data Model~2.0に準拠するよう拡張し、
    より広範な分散型アイデンティティエコシステムとの相互運用性、
    および政府発行のデジタルアイデンティティシステムとの潜在的な統合を可能にする。
    SBTベースのキャリア記録を国家デジタルアイデンティティ基盤
    （例：日本のマイナンバー制度）と連携できれば、
    フリーランスアニメーターは制度的に認知された職業活動の証明を提示できるようになり、
    記録のない独立労働者から、その専門性が銀行、入管当局、
    公的機関に読み取れる認証された専門家へと地位が変換される。

  \item \textbf{クロスインダストリー展開：}
    隣接するクリエイティブ産業（ビデオゲーム、映画VFX、音楽制作）に
    フレームワークを適応させ、一般化可能性を検証し、
    ドメイン固有の要件を特定する。

  \item \textbf{オンチェーン分析：}
    匿名化されたSBTデータを活用して業界の労働力動態を研究する集計分析ツールを開発する
    ---役割分布、プロジェクト期間のパターン、キャリアの軌跡---
    これにより業界政策にデータドリブンな知見を提供する。

  \item \textbf{レピュテーションシステムとの統合：}
    蓄積されたSBTがどのようにレピュテーションスコアに反映され、
    スタジオが信頼できるフリーランサーを特定する助けとなるかを探る。
    その際、ゼロ知識証明（zero-knowledge proofs）を通じてアニメーターのプライバシーを保護する
    （例：「このアニメーターは5作品以上にわたり100カット以上を完了した」
    という証明を、具体的なプロジェクトを明かすことなく行う）。

  \item \textbf{貢献ベースの収益分配：}
    アニメ産業における既存のブロックチェーンイニシアティブ
    （例：Animechain、Animecoin）はIPレベルの収益分配を模索してきたが、
    これらのシステムは検証可能な貢献データではなく手動で決定された配分比率に依拠している。
    AnimatorSBTのオンチェーン貢献記録は、
    スマートコントラクトベースの収益分配メカニズムの入力層として機能し得る。
    これにより、検証済みの作業に比例した自動ロイヤリティ配分
    （例：グッズやストリーミングからの二次ライセンス収益）が可能となる。
    この拡張は、アニメクリエイターが下流のIP収益化から補償を受けられないという
    長年の問題に対処するものであり、その労働こそがその価値の基盤であるにもかかわらずである。
\end{enumerate}

AnimatorSBTは、より透明で公平なアニメ制作エコシステムに向けた一歩を示すものである
---すべてのクリエイティブな貢献が記録され、検証可能であり、
それを行った本人が所有するエコシステムである。
アニメ産業がグローバルに拡大し続ける中、
この成長の基盤にいる個々のクリエイターを保護し力を与えるツールは、
単に価値があるだけでなく不可欠なものとなる。
